\section{Feature Engineering}
\label{sec:features}

In the sibling BPR and SVD projects, the feature layer produces auxiliary
signal: confidence weights, popularity priors, content features consulted
only as a cold-start fallback. Here the feature layer \emph{is} the model's
representation -- there is no separate identity-indexed embedding table that
the features merely augment. This changes what correctness means for the
pipeline: a bug in a feature transform is not a secondary signal degrading
gracefully, it is a bug in the thing being factorized.

\subsection{Declarative specification}

\texttt{conf/features.yaml} declares every feature once, consumed
identically by training and by the serving path
(\texttt{tests/test\_features\_parity.py} asserts the two never drift apart).
Each declaration carries a \emph{fit scope}:

\begin{itemize}
  \item \textbf{catalogue} -- known the moment a product is listed (brand,
        category, price bucket, ingredient tags). Always available, cold
        item or not.
  \item \textbf{profile} -- a static, self-reported user attribute (skin
        type, skin tone). Not a statistic, so not subject to the train-only
        rule; a brand-new user supplies it at signup.
  \item \textbf{train\_only} -- a vocabulary or statistic restricted to the
        training partition. Both identity features (\texttt{user\_identity},
        \texttt{item\_identity}) use this scope: the vocabulary of ID tokens
        literally does not contain an entry for a withheld entity, which is
        what makes cold-start an operationalized property of the pipeline
        rather than an assertion about it.
\end{itemize}

Table~\ref{tab:features} lists every declared feature with its family, kind
and scope; the per-feature rationale (Table~\ref{tab:feature-rationale})
documents why each one was declared.

% Generated by scripts/export_feature_tables.py -- do not edit by hand.
\begin{table}[t]
  \centering
  \caption{Declared features. ``Scope'' is the fit boundary: catalogue
  and profile features are known a priori; train\_only features --
  including both identity features -- are restricted to the training
  partition, which is what makes cold users/items genuinely lack an ID
  token rather than merely an unseen one.}
  \label{tab:features}
  \begin{tabular}{llllc}
    \toprule
    Feature & Family & Kind & Scope & PIT \\
    \midrule
    \texttt{user\_skin\_type} & user & categorical & profile & -- \\
    \texttt{user\_skin\_tone} & user & categorical & profile & -- \\
    \texttt{user\_identity} & user & identity & train\_only & -- \\
    \texttt{item\_brand} & item & categorical & catalogue & -- \\
    \texttt{item\_category} & item & categorical & catalogue & -- \\
    \texttt{item\_price\_bucket} & item & categorical & catalogue & -- \\
    \texttt{item\_ingredient\_tags} & item & multi\_categorical & catalogue & -- \\
    \texttt{item\_identity} & item & identity & train\_only & -- \\
    \texttt{user\_review\_count\_bucket} & interaction & categorical & train\_only & \checkmark \\
    \texttt{item\_popularity\_bucket} & interaction & categorical & train\_only & -- \\
    \texttt{item\_catalogue\_context} & context & categorical & train\_only & -- \\
    \bottomrule
  \end{tabular}
\end{table}


\begin{table}[h]
  \centering
  \caption{Per-feature rationale, generated from the \texttt{rationale}
  field every entry in \texttt{conf/features.yaml} is required to carry
  (\texttt{test\_features\_schema.py} rejects a feature without one).}
  \label{tab:feature-rationale}
  % Generated by scripts/export_feature_tables.py -- do not edit by hand.
\begin{description}
  \item[\texttt{user\_skin\_type}] Core demographic attribute driving taste. Known at signup, so it is available for every user including ones with zero training history -- the attribute a brand-new user supplies is exactly what lets the hybrid model place them near users with similar skin.
  \item[\texttt{user\_skin\_tone}] Secondary demographic attribute (shade matching for base/coverage products). Same availability argument as skin type.
  \item[\texttt{user\_identity}] Per-user ID feature, LightFM-style: lets the model retain personalised signal for warm users without it being the *only* signal. The vocabulary is restricted to users seen in training, so a cold user genuinely has no ID token available -- this is the feature that operationalizes "cold start" rather than merely asserting it.
  \item[\texttt{item\_brand}] Brand loyalty is a strong prior independent of any single product's review history, and is known the moment a product is catalogued.
  \item[\texttt{item\_category}] Coarse product type (serum vs. cleanser vs. sunscreen). Known at catalogue time, safe for cold items by construction.
  \item[\texttt{item\_price\_bucket}] Discretized price tier. Budget shoppers and prestige shoppers cluster on this signal regardless of brand or category.
  \item[\texttt{item\_ingredient\_tags}] The content signal that actually explains taste: a Dry-skin user's affinity is for Hyaluronic-Acid / Ceramide products specifically, not for "Moisturizers" as a category. This is what lets the model recommend a brand-new product that happens to share tags with liked products.
  \item[\texttt{item\_identity}] Mirrors user\_identity: vocabulary restricted to items seen in training, so a cold item has no ID token and must be scored through its brand/category/ingredient tokens alone.
  \item[\texttt{user\_review\_count\_bucket}] Confidence weighting: a user with 40 training reviews has a much more reliable taste signal than one with 3. Computed with a strict `<` cutoff against the user's own interaction\_order so it can never see the very review it is being used to help predict.
  \item[\texttt{item\_popularity\_bucket}] Popularity prior used by the most-popular baseline and as part of the cold-start fallback ordering. Fitted on the training partition only, so held-out interactions never inflate a product's apparent popularity.
  \item[\texttt{item\_catalogue\_context}] Distinguishes an "established" item (enough training reviews to trust its popularity/ingredient signal) from a "new\_arrival" (sparse or zero training reviews). Cold items always fall back to this value, documenting that they are indistinguishable from genuinely new catalog entries from the model's point of view.
\end{description}

\end{table}

\subsection{Point-in-time correctness}

Aggregate features derived from the interaction log (e.g.\
\texttt{user\_review\_count\_bucket}) accept a strict cutoff and filter with
\texttt{<}, never \texttt{<=}, against each user's own
\texttt{interaction\_order}. \texttt{tests/test\_features\_leakage.py}
operationalizes this directly: it injects a review positioned after the
cutoff for an existing user and asserts that no feature value for that user
changes, and separately asserts that a review positioned \emph{exactly at}
the cutoff is excluded (the boundary is exclusive, not inclusive). Item-level
aggregates (\texttt{item\_popularity\_bucket}) are instead protected by a
\texttt{fit\_on: train\_only} contract, verified by asserting the statistic
computed on the training partition differs from the same statistic computed
on the full dataset -- if they matched, the held-out interactions would
already be baked into a value the model is scored on.

\subsection{A deliberate exclusion from the retrieval representation}

\texttt{item\_popularity\_bucket} and \texttt{item\_catalogue\_context} are
declared, registered and leakage-tested like every other feature, but are
not included in the item representation the hybrid model actually
factorizes. The reason is structural rather than a matter of taste: a
withheld (cold) item has, by construction, exactly zero training
interactions, so both aggregates collapse every cold item onto the same
``unpopular'' / ``new\_arrival'' token. Including them in $\qvec_i$ would
teach the model to systematically suppress cold items through a learned
global popularity prior at precisely the moment we want to test whether it
can still retrieve them through shared brand/category/ingredient structure.
The ID-only and popularity baselines already encode an item-popularity prior
through their own bias terms -- a prior that is, for the same structural
reason, unavailable for a cold item -- so nothing is lost from the
comparison by leaving these two aggregates out of the hybrid model's own
representation; they remain fully part of the feature store, the schema and
leakage tests, and the feature table above.

\subsection{Testing surface}

Three test files operationalize the three properties above:
\texttt{test\_features\_schema.py} validates \texttt{conf/features.yaml}
against the registry (every declared transform is implemented and vice
versa, every identity feature is declared \texttt{train\_only}, every
feature documents a rationale); \texttt{test\_features\_leakage.py}
implements the point-in-time and fit-scope checks described above, plus a
label-shuffling sanity check that destroys the user--item association and
asserts ranking quality collapses to the random floor;
\texttt{test\_features\_parity.py} asserts that
\texttt{recompute\_user\_tokens}/\texttt{recompute\_item\_tokens} -- the
functions the serving path calls for an entity outside the last training
run -- reproduce exactly the tokens the batch-built feature store used
during training, for every warm user/item sampled.
